\section{Justificación del Nombre del Equipo y del Sistema}
\label{sec:justificacion_nombre}

\subsection{Origen Etimológico e Histórico}
El vocablo \textit{artesa} (del latín vulgar \textit{artēsa}, y con raíces en la tradición panadera prerromana) designa originariamente al recipiente cuadrilongo de madera, de paredes inclinadas y fondo angosto, empleado tradicionalmente por maestros panaderos para realizar el amasado manual de la harina, el agua, la sal y la masa madre. 

En la cultura gastronómica tradicional de El Salvador y de Hispanoamérica, la artesa no es simplemente un utensilio de almacenamiento; es el espacio físico de transformación donde materias primas aisladas pierden su estado individual para constituir un producto cohesionado, homogéneo y con valor agregado.

\subsection{Metáfora Sistémica y Arquitectónica}
La elección de \textbf{artesa} como denominación del sistema y del equipo de ingeniería de software responde a una analogía directa con los principios de diseño de un sistema de Planificación de Recursos Empresariales (ERP):

\begin{enumerate}
    \item \textbf{Convergencia de elementos heterogéneos:} En la panadería física, verter los ingredientes en recipientes dispersos impide el leudado correcto; análogamente, en la gestión del negocio gastronómico, dispersar la información en hojas de cálculo, cuadernos y terminales desconectadas destruye la trazabilidad. La artesa digital es la base de datos relacional y transaccional donde confluyen compras, recetas, movimientos de stock, comandas y cortes de caja.
    \item \textbf{Respeto por el proceso artesanal:} A diferencia de los grandes sistemas industriales genéricos que imponen flujos rígidos diseñados para manufactura pesada, \textbf{artesa} está modelado respetando las peculiaridades de la cocina de especialidad: productos elaborados mediante recetas compuestas, insumos medidos en fracciones (gramos, mililitros, unidades), consumo diferido por comanda y atención ágil en mostrador.
    \item \textbf{Robustez e Integridad Estructural:} Tradicionalmente tallada en una sola pieza de madera maciza para resistir la fuerza mecánica del amasado sin fracturarse ni contaminar la masa, nuestra plataforma refleja dicha solidez en su arquitectura: transacciones con aislamiento estricto en PostgreSQL, funciones atómicas y validaciones determinísticas que impiden inconsistencias contables o financieras.
\end{enumerate}

\subsection{Filosofía e Identidad del Equipo}
El equipo de desarrollo asume la filosofía del artesano: dedicación al detalle, dominio de las herramientas fundamentales (estándares web puros, SQL transaccional y buenas prácticas de ingeniería) y rechazo de la complejidad accidental. La identidad del equipo se resume en proveer a los empresarios gastronómicos locales una herramienta tecnológica de nivel empresarial que dignifique su labor y proteja la rentabilidad de su negocio.
